\section{Absolute Feuchte mit Berücksichtigung des Drucks}
KI generierter Code für die absolute Feuchte unter Berücksichtigung des Drucks:
\begin{minted}[
    bgcolor=gray!10,   % Grauer Hintergrundkasten
    fontsize=\small,   % Angenehme Schriftgröße
    breaklines         % Automatischer Zeilenumbruch
]{python}
def absolute_feuchte(temp_c, rf_prozent, druck_pa):
    e_sat = 611.2 * round(2.718281828459045 ** ((17.62 * temp_c) / (temp_c + 243.12)), 10)
    # Using numpy exp or math exp for precision
    import math
    e_sat = 611.2 * math.exp((17.62 * temp_c) / (temp_c + 243.12))
    e = (rf_prozent / 100.0) * e_sat
    af = 622.0 * e / (druck_pa - e)
    return af

af_normal = absolute_feuchte(20, 85, 101325)
af_vakuum = absolute_feuchte(20, 85, 91325)
print(f"Normal: {af_normal:.2f} g/kg")
print(f"Vakuum: {af_vakuum:.2f} g/kg")
\end{minted}
Daraus ergibt sich:
\begin{equation}
    \text{aF}(\vartheta, rF, p) = 622 \cdot \frac{\frac{rF}{100} \cdot 611{,}2 \cdot e^{\frac{17{,}62 \cdot \vartheta}{\vartheta + 243{,}12}}}{p - \left(\frac{rF}{100} \cdot 611{,}2 \cdot e^{\frac{17{,}62 \cdot \vartheta}{\vartheta + 243{,}12}}\right)} \quad \left[\frac{\text{g}}{\text{kg}}\right]
\end{equation}
Mit  
\begin{align}
e_{\text{sat}}(\vartheta) &= 611{,}2 \cdot e^{\frac{17{,}62 \cdot \vartheta}{\vartheta + 243{,}12}} \quad [\text{Pa}] \\[4ex]
e(\vartheta, rF) &= \frac{rF}{100} \cdot e_{\text{sat}}(\vartheta) \quad [\text{Pa}]
\end{align}
folgt daraus die kompakte Form der absoluten Feuchte unter Berücksichtigung des Drucks:
\begin{align}
\text{aF}(\vartheta, rF, p) &= 622 \cdot \frac{e(\vartheta, rF)}{p - e(\vartheta, rF)} \quad \left[\frac{\text{g}}{\text{kg}}\right]
\end{align}
wobei
\begin{align*}
\vartheta &= \text{Temperatur in } ^\circ\text{C} \\[3ex]
rF &= \text{Relative Luftfeuchtigkeit in } \% \, (0 \dots 100) \\[3ex]
p &= \text{Tatsächlicher Gesamtdruck (Betriebsdruck) in Pa} \\[3ex]
e_{\text{sat}} &= \text{Sättigungsdampfdruck in Pa} \\[3ex]
e &= \text{Tatsächlicher Wasserdampfpartialdruck in Pa}
\end{align*}
